% ============================================================================
% Fichier  : partie1/P01_C03_engagement.tex
% Titre    : Règles d'Engagement et Éthique de la Mission
% Auteur   : MINKA MI NGUIDJOÏ Thierry Emmanuel
% Version  : 1.0 -- novembre 2026
% ============================================================================
% Sources : NIST SP 800-115 (2008), planification ; PTES (pré-engagement) ;
% codes de déontologie (ISC2, chapitre éthique de SEC3266).
% ============================================================================

\chapter{Règles d'Engagement et Éthique de la Mission}
\label{chap:engagement}
\chaptermark{Règles d'engagement}

\epigraph{%
    « Ce qui n'est pas écrit dans les règles d'engagement\\
    est interdit. »%
}{Principe du cours}

\niveauChapitre{Fondamental}{Chapitres~\ref{chap:ethique}
et~\ref{chap:droit}.}

\begin{boxobjectifs}
\begin{itemize}
    \item Rédiger les règles d'engagement d'un test.
    \item Définir un périmètre vérifiable, adresse par adresse.
    \item Organiser la communication et l'arrêt d'urgence.
    \item Traiter les données rencontrées et les découvertes fortuites.
    \item Appliquer les principes déontologiques du testeur.
\end{itemize}
\end{boxobjectifs}

\bigskip

% ============================================================================
\section*{Question directrice : que faire quand on trouve, au milieu d'un test, une trace d'attaque réelle~?}
% ============================================================================

Le 27~octobre, Mme~NGO~MBOCK découvre sur un serveur de Garoua un compte
« svc\_maj2 » créé le 4~mars~: la trace de l'attaquant de l'Opération
GAROUA, oubliée lors du nettoyage. Ce n'est pas l'objet du test. Que
faire~? Les règles d'engagement le disent, si elles ont été bien écrites.

\medskip
\textbf{Plan du chapitre.}
\begin{enumerate}
    \item Les règles d'engagement (section~\ref{sec:c03-regles}).
    \item Le périmètre (section~\ref{sec:c03-perimetre}).
    \item Communiquer et s'arrêter (section~\ref{sec:c03-communiquer}).
    \item Données et découvertes fortuites (section~\ref{sec:c03-donnees}).
    \item La déontologie du testeur (section~\ref{sec:c03-deontologie}).
\end{enumerate}

\separateur

% ============================================================================
\section{Les Règles d'Engagement}
\label{sec:c03-regles}
% ============================================================================

\begin{boxdef}
\termecle{Règles d'engagement}~: document qui précise, pour une mission de
test, ce qui est autorisé, comment, quand, par qui, et ce qu'il faut faire
en cas d'incident. Il complète la lettre d'autorisation et la rend
opérationnelle.
\end{boxdef}

\begin{tableaucours}{L{3.4cm}L{8.8cm}}{%
    \tableauHeader{Rubrique} &
    \tableauHeader{Contenu}
}
Objectifs & Questions auxquelles le test répond, drapeaux à atteindre \\
Périmètre & Adresses, domaines, applications, comptes~; exclusions \\
Méthodes & Autorisées (balayage, exploitation contrôlée, hameçonnage
simulé) et interdites (déni de service, destruction) \\
Fenêtre & Dates, heures, fuseau horaire~; tests de nuit sur accord \\
Équipe & Noms des testeurs, adresses IP sources, postes utilisés \\
Communication & Contacts, canal chiffré, points d'avancement \\
Arrêt & Mot d'arrêt, motifs d'arrêt, reprise \\
Données & Ce qu'on peut consulter, ce qu'on ne copie jamais, destruction \\
Découvertes & Conduite à tenir face à une compromission réelle ou à une
faille critique \\
\end{tableaucours}
\vspace{-0.8em}
\begin{center}{\small\itshape Tableau~3.1. Les rubriques des règles d'engagement}\end{center}

% ============================================================================
\section{Le Périmètre}
\label{sec:c03-perimetre}
% ============================================================================

Un périmètre vérifiable se décrit en adresses et en noms précis, avec leur
propriétaire. Une plage en notation CIDR (10.20.0.0/16) couvre 65\,536
adresses~; une exclusion (10.10.50.0/24) en retire 256. Avant chaque
action, le testeur vérifie que la cible est dedans.

\begin{boxprocedure}
\textbf{Vérifier qu'une cible est dans le périmètre}
\begin{enumerate}
    \item La cible appartient-elle à une plage ou à un domaine listés~?
    \item Appartient-elle à une exclusion~? Si oui, elle est hors périmètre.
    \item Le signataire en est-il propriétaire ou exploitant~? (vérifier les
          registres d'adresses et le contrat d'hébergement)
    \item La méthode envisagée est-elle autorisée pour cette cible~?
    \item L'heure est-elle dans la fenêtre~?
\end{enumerate}
\end{boxprocedure}

\begin{boxalert}
\textbf{Violation typique.} Suivre un lien ou une redirection vers un
domaine non listé (un prestataire de paiement, un fournisseur en nuage).

\smallskip
\textbf{Règle~:} un domaine découvert pendant le test n'entre pas dans le
périmètre de lui-même. On le note et on demande un avenant écrit.
\end{boxalert}

% ============================================================================
\section{Communiquer et S'Arrêter}
\label{sec:c03-communiquer}
% ============================================================================

\begin{tableaucours}{L{4.0cm}L{8.2cm}}{%
    \tableauHeader{Situation} &
    \tableauHeader{Conduite}
}
Service ralenti ou arrêté & Arrêt immédiat, appel du contact d'urgence \\
Faille critique exploitable depuis Internet & Prévenir le client sans
attendre le rapport \\
Trace d'une compromission réelle & Arrêt du test sur la cible, appel du
contact, préservation \\
Doute sur le périmètre & Ne pas agir~; demander par écrit \\
Mot d'arrêt prononcé par le client & Arrêt de toutes les actions \\
\end{tableaucours}
\vspace{-0.8em}
\begin{center}{\small\itshape Tableau~3.2. Quand prévenir, quand s'arrêter}\end{center}

% ============================================================================
\section{Données et Découvertes Fortuites}
\label{sec:c03-donnees}
% ============================================================================

\begin{boxscenario}{Opération LOGONE. Le compte du 4 mars}
\textbf{Constatation.} Le 27~octobre à 10~h~12, Mme~NGO~MBOCK trouve,
membre des administrateurs du domaine, le compte « svc\_maj2 », créé le
4~mars à 2~h~04 (cours SEC3266, lab de journaux).

\medskip
\textbf{Conduite.} Elle s'arrête, ne touche pas au compte, note l'heure et
la preuve, appelle M.~BIKORO, qui appelle M.~KOUNOU à 10~h~25 selon les
règles d'engagement. LiZbiZ désactive le compte, ouvre un incident et fait
vérifier l'absence d'autre persistance~; le test reprend le lendemain.

\medskip
\textbf{Analyse} (confiance élevée). La découverte relève de la réponse aux
incidents, pas du test. Exploiter ce compte aurait mêlé l'activité de
l'attaquant et celle du testeur, et détruit des preuves.
\end{boxscenario}

% ============================================================================
\section{La Déontologie du Testeur}
\label{sec:c03-deontologie}
% ============================================================================

\begin{tableaucours}{L{3.6cm}L{8.6cm}}{%
    \tableauHeader{Principe} &
    \tableauHeader{Application}
}
Légalité & Rien hors mandat~; rien d'interdit par la loi \\
Minimisation & Faire le minimum nécessaire pour prouver~; ne pas aller
au-delà de l'objectif \\
Confidentialité & Résultats au seul client~; aucune publication sans
accord \\
Loyauté & Rapporter tout, y compris ses propres erreurs pendant le test \\
Compétence & Ne pas employer une technique que l'on ne maîtrise pas sur un
système de production \\
\end{tableaucours}
\vspace{-0.8em}
\begin{center}{\small\itshape Tableau~3.3. Les principes déontologiques du testeur}\end{center}

\begin{boxobjet}{T-15 · Découverte fortuite (hors constats du test)}
\textbf{Description~:} compte « svc\_maj2 » membre des administrateurs du
domaine, créé le 4~mars par l'attaquant. \textbf{Preuve~:} liste des
membres du groupe et date de création (capture du 27~octobre, 10~h~12).
\textbf{Gravité~:} critique (accès persistant d'un attaquant).
\textbf{Recommandation~:} désactiver, ouvrir un incident, rechercher toute
autre persistance. \textbf{Détection~:} la revue des comptes à privilèges
aurait dû le voir.
\end{boxobjet}

\begin{boxmetier}{Analyste SOC et équipe rouge}
\textbf{Distinguer le testeur de l'attaquant.} Les règles d'engagement
donnent au SOC les adresses sources et les fenêtres du testeur. Toute
activité hors de ces adresses ou de ces heures se traite comme une attaque
réelle.
\end{boxmetier}

\separateur

\begin{boxresume}
\textbf{Ce qu'il faut retenir de ce chapitre~:}
\begin{itemize}
    \item \textbf{Les règles d'engagement} rendent l'autorisation
          opérationnelle~: objectifs, périmètre, méthodes, fenêtre, équipe,
          communication, arrêt, données, découvertes.
    \item \textbf{Le périmètre} se vérifie adresse par adresse avant chaque
          action.
    \item \textbf{On s'arrête et on prévient} en cas d'incident, de faille
          critique ou de doute.
    \item \textbf{Une compromission réelle} découverte relève de la réponse
          aux incidents.
    \item \textbf{Légalité, minimisation, confidentialité, loyauté,
          compétence.}
\end{itemize}
\end{boxresume}

\begin{boxinfo}
\textbf{Points de contrôle de l'élève.} Je sais~: rédiger des règles
d'engagement~; vérifier qu'une cible est dans le périmètre~; dire quand
s'arrêter et qui prévenir~; traiter une découverte fortuite~; appliquer
les principes déontologiques.
\end{boxinfo}

\bigskip

\begin{boxexo}[title={\ding{45}\quad Exercice~3.1. CIDR \niveaubase}]
Combien d'adresses comptent 198.51.100.0/27 et 10.10.50.0/24~? L'adresse
198.51.100.40 est-elle dans la première plage~?
\end{boxexo}

\begin{boxexo}[title={\ding{45}\quad Exercice~3.2. Les rubriques \niveaubase}]
Quelle rubrique des règles d'engagement répond à~: « puis-je tester à
22~h~? »~; « que faire si le site tombe~? »~; « puis-je copier un fichier
pour preuve~? »
\end{boxexo}

\begin{boxexo}[title={\ding{45}\quad Exercice~3.3. Redirection \niveaubase}]
Le portail de suivi redirige vers « paiement.exemple-prestataire.cm ».
Pouvez-vous le tester~? Que faites-vous~?
\end{boxexo}

\begin{boxexo}[title={\ding{45}\quad Exercice~3.4. Faille critique \niveaumoyen}]
Vous découvrez le premier jour une faille critique exploitable depuis
Internet. Attendez-vous le rapport final~? Rédigez le message au client.
\end{boxexo}

\begin{boxexo}[title={\ding{45}\quad Exercice~3.5. Le SOC \niveaumoyen}]
Faut-il prévenir le SOC d'un test~? Comparez les avantages d'un test annoncé
et d'un test non annoncé.
\end{boxexo}

\begin{boxexo}[title={\ding{45}\quad Exercice~3.6. Minimisation \niveaumoyen}]
Donnez trois exemples de gestes qui dépassent le minimum nécessaire pour
prouver une faille.
\end{boxexo}

\begin{boxexo}[title={\ding{45}\quad Exercice~3.7. Sa propre erreur \niveaumoyen}]
Un testeur a supprimé par erreur un fichier de configuration. Que dit le
principe de loyauté~? Que doit-il faire~?
\end{boxexo}

\begin{boxexo}[title={\ding{45}\quad Exercice~3.8. Mot d'arrêt \niveaumoyen}]
Proposez la procédure d'arrêt et de reprise d'un test, avec ses
responsables.
\end{boxexo}

\begin{boxexo}[title={\ding{45}\quad Exercice~3.9. Des règles complètes \niveauexpert}]
Rédigez les règles d'engagement de l'Opération MEFOU (portail et API de
l'ANPP), selon les rubriques du Tableau~3.1.
\end{boxexo}

\begin{boxexo}[title={\ding{45}\quad Exercice~3.10. Le dilemme \niveauexpert}]
Pendant un test, vous trouvez la preuve qu'un salarié détourne des fonds.
Analysez la situation selon la méthode de résolution d'un dilemme (cours
SEC3266) et dites ce que vous faites.
\end{boxexo}

\bigskip

\begin{boxinfo}
\textbf{Pour aller plus loin.}
\begin{itemize}
    \item NIST SP~800-115, planification~\cite{nist800115}.
    \item PTES, phase de pré-engagement~\cite{ptes}.
\end{itemize}
\end{boxinfo}

\begin{boxinfo}
\textbf{Transition.} Le cadre posé, il faut une méthode pour conduire le
test~: le Chapitre~\ref{chap:methodes} présente les méthodologies et les
référentiels.
\end{boxinfo}
